\documentclass[11pt,a4paper]{article}

% Fonts 
\usepackage[utf8]{inputenc}
\usepackage[T1]{fontenc}
\usepackage{textcomp}
\usepackage{lmodern}

% Layout
\usepackage[margin=2.5cm]{geometry}
\usepackage{setspace}
\setlength{\parindent}{0pt}
\setlength{\parskip}{0.7em plus 0.2em minus 0.1em}
\onehalfspacing

% Language
\usepackage{babel}
\babelprovide[main,import]{british}

% Bibliography (APA 7th edition)
\usepackage{csquotes}
\usepackage[backend=biber,style=apa,sorting=nyt,natbib=true]{biblatex}
\DeclareLanguageMapping{british}{british-apa}
\addbibresource{bibliography.bib}

% Mathematics
\usepackage{amsmath,amssymb}

% Tables and figures
\usepackage{booktabs}
\usepackage{longtable}
\usepackage{graphicx}
\usepackage{float}
\usepackage{enumitem}

% Hyperlinks
\usepackage{xcolor}
\usepackage[colorlinks=true,linkcolor=blue!60!black,citecolor=blue!60!black,urlcolor=blue!60!black]{hyperref}

% Captions
\usepackage[font=small,labelfont=bf]{caption}

% Section formatting
\usepackage{titlesec}
\titleformat{\section}{\large\bfseries}{\thesection}{1em}{}
\titleformat{\subsection}{\normalsize\bfseries}{\thesubsection}{1em}{}
\titleformat{\subsubsection}{\normalsize\itshape}{\thesubsubsection}{1em}{}

% Other
\usepackage{array}
\usepackage{multirow}
\usepackage{tabularx}
\usepackage{pdflscape}
\usepackage{rotating}

\definecolor{sourceblue}{HTML}{D6E4F0}
\definecolor{headerblue}{HTML}{2471A3}

\newcolumntype{L}[1]{>{\raggedright\arraybackslash}p{#1}}
\newcolumntype{C}[1]{>{\centering\arraybackslash}p{#1}}

\newcommand{\metarow}[4]{\texttt{#1} & #2 & #3 & #4 \\}

\title{\textbf{Appendix 1: Dataset downloading proceedure}\\[0.5em] \large MSc Economic Policy Thesis\\Utrecht University School of Economics}
\author{Maurits de Haan\\[0.3em]
\small Supervisor: dr.\ Peter Gerbrands}
\date{March 2026}

\begin{document}
\maketitle
\tableofcontents

\section{Introduction} \label{sec:intro}

This document provides a complete overview of every dataset collected for the GMT thesis project.  All data were downloaded via the Python script \texttt{data/collect\_all\_data.py}.  

\section{OECD Country-by-Country Reporting (CbCR)} \label{sec:cbcr}

\subsection{Background}
\label{subsec:cbcr_background}

Under the OECD/G20 Base Erosion and Profit Shifting (BEPS) Action~13, large
multinational enterprise (MNE) groups with consolidated revenue exceeding
EUR~750~million are required to file Country-by-Country Reports (CbCRs)
detailing their economic activity and tax payments in each jurisdiction where
they operate.  The OECD aggregates and anonymises these reports into
jurisdiction-level statistics published in the \emph{OECD Corporate Tax
Statistics} (CTS) database.

The CbCR statistics are uniquely valuable because they capture tax-planning
behaviour at the MNE level across \emph{all} jurisdictions, including
low-tax locations typically absent from firm-level datasets.  Key measures
include pre-tax profits, taxes accrued and paid, employees, total revenue,
and assets, all disaggregated by \emph{parent jurisdiction} (REF\_AREA) and
\emph{partner jurisdiction} (COUNTERPART\_AREA).

The OECD publishes six tables of CbCR statistics. 

\subsection{Table II --- Interquartile Means by Jurisdiction}
\label{subsec:cbcr_t2}

\textbf{File:} \texttt{raw/oecd\_cbcr/cbcr\_table2\_interquartile\_means.csv}\\
\textbf{OECD dataflow:} \texttt{OECD.CTP.TPS,DSD\_CBCR@DF\_CBCRII,1.1}\\
\textbf{Rows:} 437{,}129 \quad \textbf{Size:} 43.4~MB\\
\textbf{Years:} 2016--2022 \quad \textbf{Reporting countries:} 46

Table~II reports \emph{interquartile means} of key financial aggregates for
pairs of reporting (parent) and partner jurisdictions, further disaggregated by
profit grouping (\texttt{PROFIT\_GROUPING}: positive/negative/total),
ETR bracket (\texttt{EFFECTIVE\_TAX\_RATE}), and MNE-size class
(\texttt{MNE\_SIZE}).  Using interquartile means instead of straight averages
reduces the influence of extreme outliers.

\begin{longtable}{L{2.8cm} L{5.5cm} C{1.8cm} C{2.5cm}}
    \caption{CbCR Table~II --- variable codes and descriptions}
    \label{tab:cbcr_t2}\\
    \toprule
    \textbf{Code (MEASURE)} & \textbf{Description} & \textbf{Unit} & \textbf{Coverage} \\
    \midrule
    \endfirsthead
    \multicolumn{4}{c}{\textit{(continued from previous page)}} \\
    \toprule
    \textbf{Code (MEASURE)} & \textbf{Description} & \textbf{Unit} & \textbf{Coverage} \\
    \midrule
    \endhead
    \midrule
    \multicolumn{4}{r}{\textit{(continued on next page)}} \\
    \endfoot
    \bottomrule
    \endlastfoot
    \texttt{PROFIT}         & Pre-tax profit (accounting profit before income tax) & USD (millions) & 2016--2022 \\
    \texttt{PROFIT\_ADJ}    & Profit adjusted for intra-group dividends & USD (millions) & 2016--2022 \\
    \texttt{TAX\_ACCRUED}   & Income tax accrued in the current period & USD (millions) & 2016--2022 \\
    \texttt{TAX\_PAID}      & Income tax paid (cash basis, including withholding) & USD (millions) & 2016--2022 \\
    \texttt{TOT\_REV}       & Total revenue (related and unrelated party) & USD (millions) & 2016--2022 \\
    \texttt{UPR}            & Unrelated-party revenue & USD (millions) & 2016--2022 \\
    \texttt{RPR}            & Related-party revenue & USD (millions) & 2016--2022 \\
    \texttt{EARNINGS}       & Employees' earnings (remuneration) & USD (millions) & 2016--2022 \\
    \texttt{EMPLOYEES}      & Number of employees & Persons & 2016--2022 \\
    \texttt{ASSETS}         & Total assets (book value) & USD (millions) & 2016--2022 \\
    \texttt{STATED\_CAPITAL}& Stated (paid-in) capital & USD (millions) & 2016--2022 \\
    \texttt{CBCR\_COUNT}    & Number of CbCR filers with activity in this jurisdiction pair & Entities & 2016--2022 \\
    \texttt{ENTITIES\_COUNT}& Total number of constituent entities & Entities & 2016--2022 \\
    \texttt{SUBGROUPS\_COUNT}& Number of MNE sub-groups & Entities & 2016--2022 \\
\end{longtable}

\noindent\textbf{Dimensions:}
\begin{itemize}[noitemsep]
    \item \texttt{STATISTICAL\_OPERATION}: \texttt{IQMEAN} (interquartile mean);
    also \texttt{IQSUM} (sum over IQ range).
    \item \texttt{AGGREGATION\_TYPE}: \texttt{TOTAL}, \texttt{TRSUB} (by
    subsidiary transfer), \texttt{TRMNE} (by MNE group).
    \item \texttt{PROFIT\_GROUPING}: \texttt{PANELAIR} (positive profits),
    \texttt{PANELAII} (negative profits), \texttt{TOTAL}.
    \item \texttt{EFFECTIVE\_TAX\_RATE}: ETR bucket thresholds
    (e.g.\ \texttt{\_Z} = all; \texttt{ETR\_BELOW15} = below 15~\%).
    \item \texttt{MNE\_SIZE}: total turnover thresholds for the MNE group.
\end{itemize}

\subsection{Table IV --- Aggregate Totals by ETR Bracket and Activity Type}
\label{subsec:cbcr_t4}

\textbf{File:} \texttt{raw/oecd\_cbcr/cbcr\_table4\_etr\_aggregates\_alt.csv}\\
\textbf{OECD dataflow:} \texttt{OECD.CTP.TPS,DSD\_CBCR@DF\_CBCRIV,1.1}\\
\textbf{Rows:} 101{,}399 \quad \textbf{Size:} 10.1~MB\\
\textbf{Years:} 2016--2022 \quad \textbf{Reporting countries:} 49

Table~IV provides aggregate totals (not interquartile-restricted) of the same
financial measures as Table~II, but separately by \emph{reported activity type}
of constituent entities in each jurisdiction.  This allows researchers to
identify whether profits in a low-tax jurisdiction are associated with holding,
IP, or genuine production activity.

The variables in Table~IV are identical to those in Table~II
(see Table~\ref{tab:cbcr_t2}) plus the following \textbf{activity type codes}:

\begin{longtable}{L{2.8cm} L{8.5cm}}
    \caption{CbCR Table~IV --- additional activity-type codes}
    \label{tab:cbcr_t4_act}\\
    \toprule
    \textbf{Code (MEASURE)} & \textbf{Activity description} \\
    \midrule
    \endfirsthead
    \multicolumn{2}{c}{\textit{(continued)}} \\
    \toprule
    \textbf{Code} & \textbf{Description} \\
    \midrule
    \endhead
    \bottomrule
    \endlastfoot
    \texttt{ACT\_ADMIN}   & Administrative, management, or support services \\
    \texttt{ACT\_DOR}     & Dormant entities \\
    \texttt{ACT\_FIN}     & Regulated financial services (banking, insurance) \\
    \texttt{ACT\_HOLDING} & Holding or managing shares or other equity instruments \\
    \texttt{ACT\_IGF}     & Internal group financing \\
    \texttt{ACT\_INS}     & Insurance and reinsurance \\
    \texttt{ACT\_IP}      & Holding, managing, or exploiting intellectual property \\
    \texttt{ACT\_MAN}     & Manufacturing or production \\
    \texttt{ACT\_OTHER}   & Other activities \\
    \texttt{ACT\_PUR}     & Purchasing or procurement \\
    \texttt{ACT\_RD}      & Research and development \\
    \texttt{ACT\_SALES}   & Sales, marketing, or distribution \\
    \texttt{ACT\_SER}     & Services provided to third parties \\
\end{longtable}

\subsection{Table V --- MNE-Size Distribution}
\label{subsec:cbcr_t5}

\textbf{File:} \texttt{raw/oecd\_cbcr/cbcr\_table5\_mne\_size\_distribution.csv}\\
\textbf{OECD dataflow:} \texttt{OECD.CTP.TPS,DSD\_CBCR@DF\_CBCRV,1.1}\\
\textbf{Rows:} 4{,}195 \quad \textbf{Size:} 0.4~MB\\
\textbf{Years:} 2016--2022 \quad \textbf{Reporting countries:} 40

Table~V contains distributional statistics (10th and 90th percentile values) for
MNE groups by the size of the group's ultimate parent entity, providing
information on the concentration of economic activity.

\begin{longtable}{L{2.8cm} L{5.5cm} C{1.8cm} C{2.5cm}}
    \caption{CbCR Table~V --- variable codes}
    \label{tab:cbcr_t5}\\
    \toprule
    \textbf{Code (MEASURE)} & \textbf{Description} & \textbf{Unit} & \textbf{Stat.\ op.}\\
    \midrule
    \endfirsthead
    \toprule
    \textbf{Code} & \textbf{Description} & \textbf{Unit} & \textbf{Stat.\ op.}\\
    \midrule
    \endhead
    \bottomrule
    \endlastfoot
    \texttt{ASSETS}    & Total assets of reporting MNE group & USD & P10/P90 percentiles \\
    \texttt{EMPLOYEES} & Total employees of reporting MNE group & Persons & P10/P90 percentiles \\
    \texttt{UPR}       & Unrelated-party revenue of reporting MNE group & USD & P10/P90 percentiles \\
\end{longtable}

\subsection{Table VI --- Aggregates by MNE-Group Size}
\label{subsec:cbcr_t6}

\textbf{File:} \texttt{raw/oecd\_cbcr/cbcr\_table6\_mne\_size\_aggregates.csv}\\
\textbf{OECD dataflow:} \texttt{OECD.CTP.TPS,DSD\_CBCR@DF\_CBCRVI,1.1}\\
\textbf{Rows:} 105{,}753 \quad \textbf{Size:} 10.7~MB\\
\textbf{Years:} 2022 only \quad \textbf{Reporting countries:} 39

Table~VI disaggregates the same financial measures and activity types as
Table~IV by size class of the MNE group (\texttt{MNE\_SIZE}), allowing
analysis of how profit-shifting patterns vary with the size of the multinational.
Currently only 2022 data are available.

\begin{longtable}{L{3.2cm} L{8.0cm}}
    \caption{CbCR Table~VI --- MNE size classifications}
    \label{tab:cbcr_t6_size}\\
    \toprule
    \textbf{\texttt{MNE\_SIZE} code} & \textbf{Threshold definition} \\
    \midrule
    \endfirsthead
    \toprule
    \textbf{Code} & \textbf{Definition} \\
    \midrule
    \endhead
    \bottomrule
    \endlastfoot
    \texttt{SIZE\_LESS\_1B}     & Group consolidated revenue $<$ USD~1~billion \\
    \texttt{SIZE\_1B\_10B}      & USD~1~billion $\leq$ revenue $<$ USD~10~billion \\
    \texttt{SIZE\_10B\_PLUS} / \texttt{SIZE\_GREATER\_10} & Revenue $\geq$ USD~10~billion \\
    \texttt{\_Z}                 & All size groups combined \\
\end{longtable}

\noindent\textbf{Coverage note:} CbCR data reflect \emph{voluntary participation} by
jurisdictions that have enacted CbCR legislation.  As a result, some
low-tax jurisdictions (e.g.\ Cayman Islands \texttt{CYM}) appear only as partner
(host) jurisdictions, not as reporting (parent) jurisdictions.  Figures are
also subject to minimum-filing-threshold filters and may be suppressed for
confidentiality where too few filers are present.


\section{OECD Statutory Corporate Income Tax Rates}
\label{sec:cit}

\subsection{Background}
\label{subsec:cit_background}

The OECD \emph{Tax Database} provides harmonised time series of statutory
corporate income tax (CIT) rates, compiled from legislation and administrative
sources.  The statutory rate is the headline rate applied to corporate profits
before allowances and credits; it differs from the effective rates documented
in Section~\ref{sec:etr}.  Statutory rates are the primary instrument of tax
competition because they are the publicly visible signal to potential investors.

Two datasets are collected: (1)~the main CIT rate series covering a broad set
of countries including both central/federal and combined (central + sub-central)
rates; and (2)~the combined statutory tax rate on dividend income, which
captures double-taxation of equity returns.

\subsection{Statutory CIT Rates}
\label{subsec:cit_rates}

\textbf{File:} \texttt{raw/oecd\_tax\_rates/oecd\_statutory\_cit\_rates.csv}\\
\textbf{OECD dataflow:} \texttt{OECD.CTP.TPS,DSD\_TAX\_CIT@DF\_CIT,2.0}\\
\textbf{Rows:} 13{,}273 \quad \textbf{Size:} 1.2~MB\\
\textbf{Years:} 2000--2025 \quad \textbf{Countries:} 144

\begin{longtable}{L{2.8cm} L{6.0cm} C{1.8cm} C{2.5cm}}
    \caption{Statutory CIT rates --- variable codes}
    \label{tab:cit_rates}\\
    \toprule
    \textbf{Code (MEASURE)} & \textbf{Description} & \textbf{Unit} & \textbf{Coverage} \\
    \midrule
    \endfirsthead
    \toprule
    \textbf{Code} & \textbf{Description} & \textbf{Unit} & \textbf{Coverage} \\
    \midrule
    \endhead
    \bottomrule
    \endlastfoot
    \texttt{CIT}       & Combined (central + sub-central) statutory CIT rate & \% of income & 2000--2025 \\
    \texttt{CIT\_C}    & Central government CIT rate only & \% of income & 2000--2025 \\
    \texttt{CIT\_SCG}  & Sub-central government CIT rate (e.g.\ canton, state) & \% of income & 2000--2025 \\
    \texttt{CIT\_EX\_ST} & Surtax or extra statutory charges & \% of income & 2000--2025 \\
    \texttt{CIT\_LSN}  & Local surtax / net wealth tax equivalent included in CIT & \% of income & 2000--2025 \\
\end{longtable}

\noindent\textbf{Dimension notes:}
\begin{itemize}[noitemsep]
    \item \texttt{TARGETING}: \texttt{ST} = standard rate; \texttt{RED} = reduced rate;
          \texttt{INC} = increased rate.
    \item \texttt{SECTOR}: \texttt{S1311} = general government; \texttt{\_Z} = all.
    \item \texttt{RATE\_STRUCTURE}: flat, progressive, etc.
    \item \texttt{DIV\_TREATMENT}: how intra-group dividends are treated.
\end{itemize}

\noindent\textbf{Coverage note:} The 144~countries include all OECD members plus a large
number of non-OECD jurisdictions tracked by the OECD CTS programme, covering
tax havens, Caribbean territories, and major emerging markets.  For analytical
purposes the standard combined rate (\texttt{CIT}, \texttt{TARGETING=ST}) is
typically the relevant measure.

\subsection{Combined Statutory Tax Rates on Dividend Income}
\label{subsec:dividend_rates}

\textbf{File:} \texttt{raw/oecd\_tax\_rates/oecd\_combined\_dividend\_rates.csv}\\
\textbf{OECD dataflow:} \texttt{OECD.CTP.TPS,DSD\_TAX\_CIT@DF\_CIT\_DIVD\_INCOME,1.0}\\
\textbf{Rows:} 7{,}915 \quad \textbf{Size:} 0.8~MB\\
\textbf{Years:} 2000--2025 \quad \textbf{Countries:} 38 (OECD members)

These data capture the combined burden of corporate, personal, and withholding
taxes on dividend income, important for analysing the overall tax wedge on
equity investment.

\begin{longtable}{L{2.8cm} L{6.5cm} C{1.8cm} C{2.5cm}}
    \caption{Combined dividend tax rates --- variable codes}
    \label{tab:div_rates}\\
    \toprule
    \textbf{Code (MEASURE)} & \textbf{Description} & \textbf{Unit} & \textbf{Coverage} \\
    \midrule
    \endfirsthead
    \toprule
    \textbf{Code} & \textbf{Description} & \textbf{Unit} & \textbf{Coverage} \\
    \midrule
    \endhead
    \bottomrule
    \endlastfoot
    \texttt{CITPITCIT}  & Integrated CIT + PIT rate on grossed-up dividend (classical system) & \% & 2000--2025 \\
    \texttt{CPITCIT}    & Combined CIT + PIT rate (with integration relief) & \% & 2000--2025 \\
    \texttt{CIT\_DP}    & CIT rate applied to dividend payment & \% & 2000--2025 \\
    \texttt{DP}         & Dividend payment (per unit of pre-tax profit) & ratio & 2000--2025 \\
    \texttt{FWHT}       & Final withholding tax rate on dividends & \% & 2000--2025 \\
    \texttt{GUD}        & Grossed-up dividend (per unit of after-tax profit) & ratio & 2000--2025 \\
    \texttt{IDTC}       & Imputation/dividend tax credit & \% & 2000--2025 \\
    \texttt{IR}         & Imputation rate under partial imputation systems & \% & 2000--2025 \\
    \texttt{NPT}        & Net personal tax on dividend income & \% & 2000--2025 \\
    \texttt{PITGUD}     & Personal income tax on grossed-up dividend & \% & 2000--2025 \\
    \texttt{PTDP}       & Personal tax on dividend payment & \% & 2000--2025 \\
\end{longtable}


\section{OECD Forward-Looking Effective Tax Rates}
\label{sec:etr}

\subsection{Background}
\label{subsec:etr_background}

Forward-looking effective tax rates (ETRs) measure the tax burden on a
\emph{hypothetical marginal or average investment project}, incorporating not
only the statutory rate but also all relevant reliefs: capital allowances,
depreciation rules, interest deductibility, and targeted incentives such as
IP~boxes and R\&D tax provisions.  The OECD \emph{Corporate Tax Statistics --
Forward-Looking Effective Tax Rates} dataset, first published in 2020, provides
two complementary measures:

\begin{itemize}[noitemsep]
    \item \textbf{EMTR} (\emph{effective marginal tax rate}): the proportional
    tax wedge between the pre-tax and after-tax rate of return on the last
    unit of a break-even (zero-NPV) investment.  This captures the distortion of
    the \emph{extensive} margin (whether to invest).
    \item \textbf{EATR} (\emph{effective average tax rate}): the ratio of the
    present value of taxes to the present value of pre-tax profit for a
    fixed-return investment.  This captures the margin relevant for \emph{location
    choice} decisions of multinational firms.
\end{itemize}

Three files are collected: a baseline (no incentives), a set that incorporates
\emph{income-based} incentives (IP~boxes, patent boxes), and a set that
incorporates \emph{expenditure-based} incentives (R\&D credits, enhanced
allowances).

\subsection{ETR Baseline}
\label{subsec:etr_baseline}

\textbf{File:} \texttt{raw/oecd\_etr/oecd\_etr\_baseline.csv}\\
\textbf{OECD dataflow:} \texttt{OECD.CTP.TPS,DSD\_ETR@DF\_ETR\_BASELINE,2.0}\\
\textbf{Rows:} 29{,}489 \quad \textbf{Size:} 3.0~MB\\
\textbf{Years:} 2017--2024 \quad \textbf{Countries:} 103

\begin{longtable}{L{2.8cm} L{6.0cm} C{2.0cm} C{2.5cm}}
    \caption{ETR baseline --- variable codes}
    \label{tab:etr_baseline}\\
    \toprule
    \textbf{Code (MEASURE)} & \textbf{Description} & \textbf{Unit} & \textbf{Coverage} \\
    \midrule
    \endfirsthead
    \toprule
    \textbf{Code} & \textbf{Description} & \textbf{Unit} & \textbf{Coverage} \\
    \midrule
    \endhead
    \bottomrule
    \endlastfoot
    \texttt{EATR}               & Effective average tax rate (no incentives) & \% of income & 2017--2024 \\
    \texttt{EMTR}               & Effective marginal tax rate (no incentives) & \% of income & 2017--2024 \\
    \texttt{PBAR}               & $\bar{p}$: pre-tax rate of return on break-even investment & \% & 2017--2024 \\
    \texttt{CAPITAL\_ALLOWANCES}& Present value of capital allowances (NPV of depreciation deductions) & \% of cost & 2017--2024 \\
\end{longtable}

\noindent\textbf{Dimension notes:}
\begin{itemize}[noitemsep]
    \item \texttt{ETR\_TAX\_BASIS}: \texttt{BASELINE} = no incentives; also
    distinguishes \texttt{CORPORATE} from \texttt{SHAREHOLDER}-level.
    \item \texttt{ETR\_TAX\_TYPE}: \texttt{CS} = consolidated CIT;
    \texttt{FIXED} = fixed-rate assumption; \texttt{AWS} = actual weighted-average.
    \item \texttt{UNIT\_MEASURE}: \texttt{PT\_INC\_TAX} = as \% of income tax;
    \texttt{PT\_P5L} = as \% of initial cost; \texttt{PT\_P5L\_INIT}.
\end{itemize}

\subsection{ETR with Income-Based Incentives (IP~Boxes)}
\label{subsec:etr_income}

\textbf{File:} \texttt{raw/oecd\_etr/oecd\_etr\_income\_incentives.csv}\\
\textbf{OECD dataflow:} \texttt{OECD.CTP.TPS,DSD\_ETR@DF\_ETR\_IBRD,2.0}\\
\textbf{Rows:} 5{,}784 \quad \textbf{Size:} 0.7~MB\\
\textbf{Years:} 2000--2024 \quad \textbf{Countries:} 55 (with active IP regimes)

\begin{longtable}{L{2.8cm} L{6.0cm} C{2.0cm} C{2.5cm}}
    \caption{ETR with income-based incentives --- variable codes}
    \label{tab:etr_income}\\
    \toprule
    \textbf{Code (MEASURE)} & \textbf{Description} & \textbf{Unit} & \textbf{Coverage} \\
    \midrule
    \endfirsthead
    \toprule
    \textbf{Code} & \textbf{Description} & \textbf{Unit} & \textbf{Coverage} \\
    \midrule
    \endhead
    \bottomrule
    \endlastfoot
    \texttt{EATR}       & EATR \emph{without} income-based incentive applied & \% & 2000--2024 \\
    \texttt{EATR\_ENH}  & EATR \emph{with} income-based incentive (e.g.\ IP~box rate) & \% & 2000--2024 \\
    \texttt{PBAR}       & Pre-tax break-even rate \emph{without} incentive & \% & 2000--2024 \\
    \texttt{PBAR\_ENH}  & Pre-tax break-even rate \emph{with} income-based incentive & \% & 2000--2024 \\
\end{longtable}

\noindent Each combination is reported per \texttt{REGIME}: \texttt{REGIME1}
through \texttt{REGIME8} label specific IP~box or patent~box regimes in a given
country-year; \texttt{\_N} indicates the standard (non-incentive) regime.

\subsection{ETR with Expenditure-Based Incentives (R\&D)}
\label{subsec:etr_expend}

\textbf{File:} \texttt{raw/oecd\_etr/oecd\_etr\_expenditure\_incentives.csv}\\
\textbf{OECD dataflow:} \texttt{OECD.CTP.TPS,DSD\_ETR@DF\_ETR\_XBRD,2.0}\\
\textbf{Rows:} 1{,}224 \quad \textbf{Size:} 0.1~MB\\
\textbf{Years:} 2019--2024 \quad \textbf{Countries:} 51

Variable codes are identical to Table~\ref{tab:etr_income}, but the
\texttt{ENH} variants incorporate expenditure-based R\&D tax credits or
enhanced capital allowances rather than preferential income rates.  This
dataset enables comparison of the incentive value of volume credits versus
IP~box regimes.


\section{Wage and Earnings Data}
\label{sec:wages}

\subsection{Background}
\label{subsec:wages_background}

Labour cost data are collected from two sources to maximise country coverage.
OECD average annual wages provide high-quality, internationally comparable
data for OECD member countries, while ILO hourly earnings extend coverage to
a wider set of developing countries.  These data serve multiple roles in the
thesis: as direct outcome variables if analysing labour-tax incidence, as
control variables in profit-shifting regressions (since wages proxy for labour
market conditions), and as inputs for constructing wage-to-profit ratios.

\subsection{OECD Average Annual Wages}
\label{subsec:oecd_wages}

\textbf{File:} \texttt{raw/oecd\_wages/oecd\_avg\_annual\_wages.csv}\\
\textbf{OECD dataflow:} \texttt{OECD.ELS.SAE,DSD\_EARNINGS@AV\_AN\_WAGE,1.0}\\
\textbf{Rows:} 3{,}924 \quad \textbf{Size:} 0.3~MB\\
\textbf{Years:} 1990--2024 \quad \textbf{Countries:} 41 (OECD members + OECD aggregate)

\begin{longtable}{L{2.4cm} L{5.5cm} C{2.2cm} C{2.5cm}}
    \caption{OECD average annual wages --- variable codes}
    \label{tab:oecd_wages}\\
    \toprule
    \textbf{Code (MEASURE)} & \textbf{Description} & \textbf{Unit} & \textbf{Coverage} \\
    \midrule
    \endfirsthead
    \toprule
    \textbf{Code} & \textbf{Description} & \textbf{Unit} & \textbf{Coverage} \\
    \midrule
    \endhead
    \bottomrule
    \endlastfoot
    \texttt{WG} & Average annual wages of full-time dependent employees & National currency; also USD PPP & 1990--2024 \\
\end{longtable}

\noindent\textbf{Dimension notes:}
\begin{itemize}[noitemsep]
    \item \texttt{UNIT\_MEASURE}: \texttt{USD\_PPP} (2024 constant PPP USD) is the
    primary comparator; national currencies also available.
    \item \texttt{PRICE\_BASE}: \texttt{Q} = current prices; \texttt{R} = constant
    (real) prices.
    \item \texttt{AGGREGATION\_OPERATION}: \texttt{MEAN}.
    \item \texttt{SEX}: total (\texttt{\_Z}), male, or female.
\end{itemize}

\subsection{OECD Hourly Earnings Index}
\label{subsec:oecd_hourly}

\textbf{File:} \texttt{raw/oecd\_wages/oecd\_hourly\_earnings.csv}\\
\textbf{OECD dataflow:} \texttt{OECD.SDD.TPS,DSD\_EAR@DF\_HOU\_EAR,1.0}\\
\textbf{Rows:} 41{,}421 \quad \textbf{Size:} 3.4~MB\\
\textbf{Years:} 1955--2026 \quad \textbf{Countries/groups:} 41

\begin{longtable}{L{2.4cm} L{6.0cm} C{2.2cm} C{2.5cm}}
    \caption{OECD hourly earnings --- variable codes}
    \label{tab:oecd_hourly}\\
    \toprule
    \textbf{Code (MEASURE)} & \textbf{Description} & \textbf{Unit} & \textbf{Coverage} \\
    \midrule
    \endfirsthead
    \toprule
    \textbf{Code} & \textbf{Description} & \textbf{Unit} & \textbf{Coverage} \\
    \midrule
    \endhead
    \bottomrule
    \endlastfoot
    \texttt{EAR} & Index of hourly earnings for employees (base year 2015 = 100) & Index (IX) & 1955--2026 \\
\end{longtable}

\noindent\textbf{Dimension notes:}
\begin{itemize}[noitemsep]
    \item \texttt{ADJUSTMENT}: \texttt{N} = not seasonally adjusted.
    \item \texttt{ACTIVITY}: \texttt{C} = total economy; sector breakdowns available.
    \item \texttt{SECTOR}: \texttt{S1} = whole economy.
    \item \texttt{FREQ}: monthly (\texttt{M}) or quarterly (\texttt{Q}).
    \item Coverage includes area aggregates: \texttt{EA19} (euro area),
    \texttt{EU27\_2020} (EU 27), \texttt{G7}, \texttt{OECD}, \texttt{OECDE}.
\end{itemize}

\noindent\textbf{Coverage note:} This series is an \emph{index} rather than a level;
for cross-country wage comparisons the OECD average annual wages dataset
(Section~\ref{subsec:oecd_wages}) is more appropriate.

\subsection{ILO Global Hourly Earnings}
\label{subsec:ilo_wages}

\textbf{File:} \texttt{raw/ilo\_wages/ilo\_hourly\_earnings.csv}\\
\textbf{ILO dataflow:} \texttt{ILO,DF\_EAR\_EHRA\_SEX\_ECO\_CUR\_NB,1.0}\\
\textbf{Rows:} 227{,}317 \quad \textbf{Size:} 50.4~MB\\
\textbf{Years:} 1990--2025 \quad \textbf{Countries:} 105

The ILO dataset extends hourly earnings data to a large number of non-OECD
countries, many of which are not covered by the OECD wage series.  Data are
compiled from national household surveys, enterprise surveys, and administrative
records reported to the ILO; methodological notes are stored in the
\texttt{NOTE\_SOURCE} and \texttt{NOTE\_INDICATOR} columns.

\begin{longtable}{L{2.8cm} L{5.5cm} C{2.2cm} C{2.5cm}}
    \caption{ILO hourly earnings --- variable codes}
    \label{tab:ilo_wages}\\
    \toprule
    \textbf{Code (MEASURE)} & \textbf{Description} & \textbf{Unit} & \textbf{Coverage} \\
    \midrule
    \endfirsthead
    \toprule
    \textbf{Code} & \textbf{Description} & \textbf{Unit} & \textbf{Coverage} \\
    \midrule
    \endhead
    \bottomrule
    \endlastfoot
    \texttt{EAR\_EHRA\_NB} & Mean nominal hourly earnings of employees & Local currency (CR); PPP USD also flagged & 1990--2025 \\
\end{longtable}

\noindent\textbf{Dimension notes:}
\begin{itemize}[noitemsep]
    \item \texttt{SEX}: \texttt{SEX\_T} = total; also male and female.
    \item \texttt{ECO}: economic sector (\texttt{ECO\_SECTOR\_TOTAL} = economy-wide).
    \item \texttt{CUR}: \texttt{CUR\_TYPE\_LCU} = local currency; \texttt{CUR\_TYPE\_PPP} = PPP.
    \item \texttt{UNIT\_MEASURE\_TYPE}: \texttt{NB} = number (level); \texttt{CR} = currency.
    \item \texttt{OBS\_STATUS}: \texttt{B} = break in series; \texttt{E} = estimated.
    \item Confidence bounds reported in \texttt{UPPER\_BOUND} / \texttt{LOWER\_BOUND} where
    available.
\end{itemize}

\noindent\textbf{Coverage note:} Coverage is uneven across countries and years; many
non-OECD observations are spaced several years apart. The OECD wage data
(Sections~\ref{subsec:oecd_wages}--\ref{subsec:oecd_hourly}) should be
preferred for OECD countries; ILO data fill the gap for non-OECD members.


\section{World Bank Macroeconomic Controls}
\label{sec:worldbank}

\subsection{Background}
\label{subsec:wb_background}

World Bank data provide three types of country-level controls commonly included
in cross-country tax regressions: GDP per capita (economic development and
market size), population (country size), and income group classification (for
sub-group and heterogeneity analyses).  Data are drawn from the World Bank
Development Indicators API for GDP and population, and from the official income
classification methodology for the classification series.

\subsection{GDP Per Capita}
\label{subsec:wb_gdp}

\textbf{File:} \texttt{raw/world\_bank/wb\_gdp\_per\_capita.csv}\\
\textbf{Indicator:} \texttt{NY.GDP.PCAP.CD} (GDP per capita, current USD)\\
\textbf{Rows:} $\approx$2{,}170 \quad \textbf{Size:} $<$0.1~MB\\
\textbf{Years:} 2015--2024 \quad \textbf{Countries/groups:} $\approx$217

\begin{longtable}{L{2.8cm} L{5.5cm} C{2.2cm} C{2.5cm}}
    \caption{World Bank GDP per capita --- variables}
    \label{tab:wb_gdp}\\
    \toprule
    \textbf{Column} & \textbf{Description} & \textbf{Unit} & \textbf{Coverage} \\
    \midrule
    \endfirsthead
    \toprule
    \textbf{Column} & \textbf{Description} & \textbf{Unit} & \textbf{Coverage} \\
    \midrule
    \endhead
    \bottomrule
    \endlastfoot
    \texttt{iso3}        & ISO~3166-1 alpha-3 country code & --- & --- \\
    \texttt{country}     & Country name & --- & --- \\
    \texttt{year}        & Year & Integer & 2015--2024 \\
    \texttt{gdp\_pc\_usd}& GDP per capita, current USD & USD & 2015--2024 \\
\end{longtable}

\noindent\textbf{Coverage note:} The file includes World Bank regional and income-group
aggregates (e.g.\ \texttt{ZH}: Africa Eastern and Southern).  These should be
filtered out when merging with country-level data.

\subsection{Population}
\label{subsec:wb_pop}

\textbf{File:} \texttt{raw/world\_bank/wb\_population.csv}\\
\textbf{Indicator:} \texttt{SP.POP.TOTL} (total population)\\
\textbf{Rows:} $\approx$2{,}170 \quad \textbf{Size:} $<$0.1~MB\\
\textbf{Years:} 2015--2024 \quad \textbf{Countries/groups:} $\approx$217

\begin{longtable}{L{2.8cm} L{5.5cm} C{2.2cm} C{2.5cm}}
    \caption{World Bank population --- variables}
    \label{tab:wb_pop}\\
    \toprule
    \textbf{Column} & \textbf{Description} & \textbf{Unit} & \textbf{Coverage} \\
    \midrule
    \endfirsthead
    \toprule
    \textbf{Column} & \textbf{Description} & \textbf{Unit} & \textbf{Coverage} \\
    \midrule
    \endhead
    \bottomrule
    \endlastfoot
    \texttt{iso3}       & ISO~3166-1 alpha-3 country code & --- & --- \\
    \texttt{country}    & Country name & --- & --- \\
    \texttt{year}       & Year & Integer & 2015--2024 \\
    \texttt{population} & Total population & Persons & 2015--2024 \\
\end{longtable}

\subsection{Income Group Classifications}
\label{subsec:wb_income}

\textbf{File:} \texttt{raw/world\_bank/wb\_income\_classifications.csv}\\
\textbf{Source:} World Bank country-level metadata API (\texttt{/v2/country})\\
\textbf{Rows:} 299 \quad \textbf{Size:} $<$0.1~MB\\
\textbf{Type:} Cross-section (current classification as of API retrieval)

\begin{longtable}{L{2.8cm} L{6.0cm} C{2.0cm}}
    \caption{World Bank income classifications --- variables}
    \label{tab:wb_income}\\
    \toprule
    \textbf{Column} & \textbf{Description} & \textbf{Coverage} \\
    \midrule
    \endfirsthead
    \toprule
    \textbf{Column} & \textbf{Description} & \textbf{Coverage} \\
    \midrule
    \endhead
    \bottomrule
    \endlastfoot
    \texttt{iso3}         & ISO~3166-1 alpha-3 code & All 299 entries \\
    \texttt{iso2}         & ISO~3166-1 alpha-2 code & All 299 entries \\
    \texttt{name}         & Country/territory name & All 299 entries \\
    \texttt{region}       & World Bank regional grouping & All 299 entries \\
    \texttt{income\_level}& Income group label (High / Upper-middle / Lower-middle / Low / Aggregates) & All 299 entries \\
    \texttt{income\_id}   & Income group code (\texttt{HIC} / \texttt{UMC} / \texttt{LMC} / \texttt{LIC} / \texttt{NA}) & All 299 entries \\
\end{longtable}

\noindent\textbf{Coverage note:} Classifications reflect the \emph{current}
World Bank fiscal-year groupings (as of download date).  For analyses that
require time-varying income classifications, the \texttt{wb\_historical\_income\_groups.xlsx}
file (also collected, stored alongside the CSVs) contains the full historical
series from the OGHIST spreadsheet.


\section{OECD BEPS Implementation Indicators}
\label{sec:beps}

\subsection{Background}
\label{subsec:beps_background}

The OECD BEPS (Base Erosion and Profit Shifting) project produced 15 Actions to
curb aggressive tax planning.  The OECD's Questionnaire on Quality and
Distribution of Data (QDD) programme tracks countries' implementation of several
key Actions: Action~2 (hybrid mismatches), Action~4 (interest limitation),
Action~5 (IP~regimes), Action~12 (mandatory disclosure), and Action~13 (CbCR
requirements), as well as Controlled Foreign Company (CFC) rules.  These
datasets provide \emph{policy indicator} variables coding whether specific laws
or rules are in place, their thresholds, and when they were introduced.

\subsection{CbCR Requirements}
\label{subsec:beps_cbcr_req}

\textbf{File:} \texttt{raw/pillar2\_beps/oecd\_cbcr\_requirements.csv}\\
\textbf{OECD dataflow:} \texttt{OECD.CTP.TPS,DSD\_QDD\_CBCREQ@DF\_QDD\_CBCREQ,1.0}\\
\textbf{Rows:} 493 \quad \textbf{Size:} $<$0.1~MB\\
\textbf{Type:} Cross-section \quad \textbf{Countries:} 119

\begin{longtable}{L{2.8cm} L{7.5cm} C{2.5cm}}
    \caption{CbCR requirements --- variable codes}
    \label{tab:cbcr_req}\\
    \toprule
    \textbf{Code (MEASURE)} & \textbf{Description} & \textbf{Coverage} \\
    \midrule
    \endfirsthead
    \toprule
    \textbf{Code} & \textbf{Description} & \textbf{Coverage} \\
    \midrule
    \endhead
    \bottomrule
    \endlastfoot
    \texttt{CBCR\_LAW}   & Whether country has enacted domestic CbCR law (\texttt{Yes/No}) & 119 countries \\
    \texttt{THRESHOLD}   & Revenue threshold for CbCR obligation (EUR millions) & 119 countries \\
    \texttt{DEADLINE}    & Filing deadline (months after fiscal year-end) & 119 countries \\
    \texttt{UPE}         & Ultimate parent entity jurisdiction indicator & 119 countries \\
    \texttt{COMMENTS}    & Qualitative notes on implementation & 119 countries \\
\end{longtable}

\subsection{CFC Rules}
\label{subsec:beps_cfc}

\textbf{File:} \texttt{raw/pillar2\_beps/oecd\_cfc\_rules.csv}\\
\textbf{OECD dataflow:} \texttt{OECD.CTP.TPS,DSD\_QDD\_CFC@DF\_QDD\_CFC,2.0}\\
\textbf{Rows:} 1{,}406 \quad \textbf{Size:} 0.4~MB\\
\textbf{Years:} 2019, 2024, 2025 \quad \textbf{Countries:} 146

CFC rules tax the passive income of foreign subsidiaries at the parent level when
the subsidiary pays little or no tax.  They are a key backstop against offshore
profit shifting.

\begin{longtable}{L{2.8cm} L{6.5cm} C{2.5cm}}
    \caption{CFC rules --- variable codes}
    \label{tab:cfc_rules}\\
    \toprule
    \textbf{Code (MEASURE)} & \textbf{Description} & \textbf{Coverage} \\
    \midrule
    \endfirsthead
    \toprule
    \textbf{Code} & \textbf{Description} & \textbf{Coverage} \\
    \midrule
    \endhead
    \bottomrule
    \endlastfoot
    \texttt{CFC}           & Whether country has CFC rules (\texttt{Yes/No}) & 146 countries \\
    \texttt{DEF}           & Definition used for CFC (control threshold, income type) & 146 countries \\
    \texttt{INC}           & Income types covered by CFC rules & 146 countries \\
    \texttt{EXE}           & Exemptions from CFC regime & 146 countries \\
    \texttt{SUB}           & Substance exclusions & 146 countries \\
    \texttt{SUB\_DESC}     & Description of substance carve-out & 146 countries \\
    \texttt{TRIGGER\_RATE} & Low-tax threshold rate triggering CFC charge (\%) & 146 countries \\
    \texttt{YEAR}          & Year of first introduction of CFC rules & 146 countries \\
\end{longtable}

\subsection{Hybrid Mismatch Rules (BEPS Action~2)}
\label{subsec:beps_hybrid}

\textbf{File:} \texttt{raw/pillar2\_beps/oecd\_hybrid\_mismatch.csv}\\
\textbf{OECD dataflow:} \texttt{OECD.CTP.TPS,DSD\_QDD\_ACT2@DF\_QDD\_ACT2,1.0}\\
\textbf{Rows:} 985 \quad \textbf{Size:} 0.1~MB\\
\textbf{Years:} 2025 \quad \textbf{Countries:} 145

Hybrid mismatch rules neutralise tax outcomes from arrangements that exploit
differences in countries' treatment of financial instruments or entities,
typically used to generate deductions in one country without corresponding
income inclusion in another.

\begin{longtable}{L{2.8cm} L{6.5cm} C{2.5cm}}
    \caption{Hybrid mismatch rules --- key variable codes}
    \label{tab:hybrid}\\
    \toprule
    \textbf{Code (MEASURE)} & \textbf{Description} & \textbf{Coverage} \\
    \midrule
    \endfirsthead
    \toprule
    \textbf{Code} & \textbf{Description} & \textbf{Coverage} \\
    \midrule
    \endhead
    \bottomrule
    \endlastfoot
    \texttt{ACT2}     & Whether Action~2 rules are in place (\texttt{Yes/No}) & 145 countries \\
    \texttt{BR\_REC1}--\texttt{BR\_REC5} & Implementation of each primary guidance branch rule (Yes/No) & 145 countries \\
    \texttt{REC1}--\texttt{REC15}  & Implementation of specific OECD recommendations (Yes/No) & 145 countries \\
    \texttt{RULE\_DESC} & Free-text description of domestic rule & 145 countries \\
    \texttt{YEAR}       & Year of legislative enactment & 145 countries \\
\end{longtable}

\subsection{Interest Limitation Rules (BEPS Action~4)}
\label{subsec:beps_interest}

\textbf{File:} \texttt{raw/pillar2\_beps/oecd\_interest\_limitation.csv}\\
\textbf{OECD dataflow:} \texttt{OECD.CTP.TPS,DSD\_QDD\_ILR@DF\_QDD\_ILR,2.0}\\
\textbf{Rows:} 3{,}633 \quad \textbf{Size:} 0.3~MB\\
\textbf{Years:} 2019, 2024, 2025 \quad \textbf{Countries:} 146

Interest limitation rules cap the amount of interest expense deductible against
taxable income, typically to a fraction of EBITDA.  They curtail debt-shifting,
a common mechanism by which MNEs reduce profits in high-tax jurisdictions.

\begin{longtable}{L{2.8cm} L{6.5cm} C{2.5cm}}
    \caption{Interest limitation rules --- variable codes}
    \label{tab:interest}\\
    \toprule
    \textbf{Code (MEASURE)} & \textbf{Description} & \textbf{Coverage} \\
    \midrule
    \endfirsthead
    \toprule
    \textbf{Code} & \textbf{Description} & \textbf{Coverage} \\
    \midrule
    \endhead
    \bottomrule
    \endlastfoot
    \texttt{ILR}         & Whether interest limitation rule is in place (\texttt{Yes/No}) & 146 countries \\
    \texttt{GR}          & The general ratio threshold used (e.g.\ 30\% of EBITDA) & 146 countries \\
    \texttt{GR\_DESC}    & Description of general ratio rule & 146 countries \\
    \texttt{TR}          & Safe harbour (de minimis) threshold (EUR) & 146 countries \\
    \texttt{TR\_DESC}    & Description of threshold & 146 countries \\
    \texttt{FAM}         & Group ratio / equity escape provision (Yes/No) & 146 countries \\
    \texttt{CF}          & Carry-forward of disallowed interest allowed (Yes/No) & 146 countries \\
    \texttt{CF\_YEARS}   & Maximum carry-forward years & 146 countries \\
    \texttt{RPD}         & Restriction on related-party payments & 146 countries \\
    \texttt{TPD}         & Transfer-pricing documentation requirement & 146 countries \\
    \texttt{IRD}         & Interest rate documentation requirement & 146 countries \\
    \texttt{DMT}         & Debt-to-equity or debt-to-asset threshold & 146 countries \\
    \texttt{LA}          & Thin capitalisation or alternative limitation approach & 146 countries \\
    \texttt{RV}          & Revenue threshold (de minimis) for rule application & 146 countries \\
    \texttt{EXCL\_PAYER} & Excluded categories of payer & 146 countries \\
    \texttt{EXCL\_PAYMENT}& Excluded categories of interest payment & 146 countries \\
    \texttt{EXCL\_REC}   & Excluded categories of recipient & 146 countries \\
    \texttt{EXCL\_OTHER} & Other exclusions & 146 countries \\
    \texttt{OTHER\_MECH} & Additional limitation mechanisms & 146 countries \\
    \texttt{RULE\_TYPE}  & Type of rule (fixed ratio, group ratio, etc.) & 146 countries \\
    \texttt{RULE\_DESC}  & General description of rule & 146 countries \\
    \texttt{YEAR}        & Year of introduction & 146 countries \\
\end{longtable}

\subsection{Intellectual Property Regimes (BEPS Action~5)}
\label{subsec:beps_ip}

\textbf{File:} \texttt{raw/pillar2\_beps/oecd\_ip\_regimes.csv}\\
\textbf{OECD dataflow:} \texttt{OECD.CTP.TPS,DSD\_QDD\_IPR@DF\_QDD\_IPR,2.1}\\
\textbf{Rows:} 2{,}282 \quad \textbf{Size:} 0.3~MB\\
\textbf{Years:} 2018--2025 \quad \textbf{Countries:} 64 (countries with active IP regime)

IP~box (patent box) regimes tax qualifying IP income at a preferential rate,
attracting mobile intangible assets.  BEPS Action~5 introduced the
\emph{modified nexus approach}, linking relief to genuine R\&D expenditure.
Each regime is identified by \texttt{REGIME} (REGIME1, REGIME2, \ldots) and
the following variables:

\begin{longtable}{L{2.8cm} L{7.0cm} C{2.5cm}}
    \caption{IP regimes --- variable codes}
    \label{tab:ip_regimes}\\
    \toprule
    \textbf{Code (MEASURE)} & \textbf{Description} & \textbf{Coverage} \\
    \midrule
    \endfirsthead
    \toprule
    \textbf{Code} & \textbf{Description} & \textbf{Coverage} \\
    \midrule
    \endhead
    \bottomrule
    \endlastfoot
    \texttt{NAME}         & Name of IP~box regime & 64 countries \\
    \texttt{RATE}         & Preferential tax rate applicable to IP income (\%) & 64 countries \\
    \texttt{TYPE}         & Type of qualifying income (royalties, capital gains, etc.) & 64 countries \\
    \texttt{ELI\_IP}      & Types of eligible intellectual property assets & 64 countries \\
    \texttt{NEXUS}        & Whether modified nexus approach has been adopted (Yes/No) & 64 countries \\
    \texttt{YEAR\_START}  & Year regime was introduced & 64 countries \\
    \texttt{YEAR\_END}    & Year regime was closed/abolished (if applicable) & 64 countries \\
\end{longtable}

\subsection{Mandatory Disclosure Rules (BEPS Action~12)}
\label{subsec:beps_mdr}

\textbf{File:} \texttt{raw/pillar2\_beps/oecd\_mandatory\_disclosure.csv}\\
\textbf{OECD dataflow:} \texttt{OECD.CTP.TPS,DSD\_QDD\_ACT12@DF\_QDD\_ACT12,1.0}\\
\textbf{Rows:} 985 \quad \textbf{Size:} 0.1~MB\\
\textbf{Years:} 2025 \quad \textbf{Countries:} 145

Mandatory disclosure rules (MDRs) require intermediaries or taxpayers to
proactively disclose potentially aggressive tax arrangements to tax authorities.

\begin{longtable}{L{2.8cm} L{7.0cm} C{2.5cm}}
    \caption{Mandatory disclosure rules --- variable codes}
    \label{tab:mdr}\\
    \toprule
    \textbf{Code (MEASURE)} & \textbf{Description} & \textbf{Coverage} \\
    \midrule
    \endfirsthead
    \toprule
    \textbf{Code} & \textbf{Description} & \textbf{Coverage} \\
    \midrule
    \endhead
    \bottomrule
    \endlastfoot
    \texttt{ACT12}       & Whether MDR is in place (\texttt{Yes/No}) & 145 countries \\
    \texttt{RULE\_DESC}  & Description of the disclosure rule & 145 countries \\
    \texttt{YEAR}        & Year of introduction & 145 countries \\
\end{longtable}


\section{Pillar Two Global Minimum Tax --- Implementation Status}
\label{sec:pillar2}

\subsection{Background}
\label{subsec:p2_background}

BEPS Pillar Two introduces a global minimum effective tax rate of 15\% for
large MNE groups (consolidated revenue $\geq$ EUR~750~million) through a set of
interlocking rules: the Income Inclusion Rule (IIR), the Undertaxed Payments
Rule (UTPR), and the Qualified Domestic Minimum Top-up Tax (QDMTT).  Countries
implementing a QDMTT can reclaim top-up tax on domestic profits before other
countries levy it via the IIR/UTPR, creating a strong incentive for adoption.

Two datasets capture Pillar Two implementation status:

\subsection{PwC Pillar Two Country Tracker}
\label{subsec:p2_pwc}

\textbf{File:} \texttt{raw/pillar2\_qdmtt/pwc\_pillar2\_tracker.csv}\\
\textbf{Source:} PricewaterhouseCoopers global Pillar Two tracker\\
\textbf{Rows:} multiple rows per country-rule \quad \textbf{Size:} $<$0.1~MB\\
\textbf{Type:} Panel (requirement $\times$ effective-date level)

\begin{longtable}{L{2.8cm} L{6.5cm} C{3.5cm}}
    \caption{PwC Pillar Two tracker --- variable codes}
    \label{tab:pwc_p2}\\
    \toprule
    \textbf{Column} & \textbf{Description} & \textbf{Values / Format} \\
    \midrule
    \endfirsthead
    \toprule
    \textbf{Column} & \textbf{Description} & \textbf{Values / Format} \\
    \midrule
    \endhead
    \bottomrule
    \endlastfoot
    \texttt{Territory}        & Country or territory name & --- \\
    \texttt{Requirement}      & Type of Pillar Two rule & \texttt{GIR}, \texttt{IIR}, \texttt{UTPR}, \texttt{QDMTT} \\
    \texttt{Requirement Group}& Grouped requirement type & E.g.\ \texttt{GIR} \\
    \texttt{P2 Start}         & Effective start date of the rule & \texttt{YYYY-MM-DD} \\
    \texttt{P2 End}           & End/expiry date (often far future) & \texttt{YYYY-MM-DD} \\
    \texttt{FY Start Mths+}   & Fiscal year offset months & Integer \\
    \texttt{Year}/\texttt{Month}/\texttt{Day} & Wildcard (\texttt{*}) or specific date components & \texttt{*} or integer \\
    \texttt{Caveat}           & Qualitative note on transition rules or exceptions & Free text \\
    \texttt{Description}      & Explanation of the rule entry & Free text \\
    \texttt{Status}           & Implementation status & e.g.\ \texttt{Enacted}, \texttt{Consultation} \\
\end{longtable}

\noindent\textbf{Requirement codes:}
\begin{itemize}[noitemsep]
    \item \texttt{GIR} / \texttt{IIR}: Income Inclusion Rule --- top-up tax applied
    at the parent-jurisdiction level on low-taxed foreign subsidiaries.
    \item \texttt{UTPR}: Undertaxed Payments Rule --- backstop allowing other
    countries to impose the top-up tax when the IIR is not applied.
    \item \texttt{QDMTT}: Qualified Domestic Minimum Top-up Tax --- domestic
    top-up tax applied before the IIR/UTPR, retaining the revenue in the source
    jurisdiction.
\end{itemize}

\noindent\textbf{Coverage note:} Only countries that have taken legislative steps appear
in the tracker; non-adopters are absent.  The tracker is updated regularly by
PwC and reflects the position as of the download date.


\section{Auxiliary and Supplementary Sources}
\label{sec:auxiliary}

\subsection{BEA US Multinational Enterprise Data}
\label{subsec:bea}

\textbf{Target files:} \texttt{raw/bea\_mne/} (Excel format, not CSV)\\
\textbf{Source:} US Bureau of Economic Analysis (BEA), Survey of US Direct
Investment Abroad\\
\textbf{Status:} Collection script configured; files not present in current
dataset (requires download from BEA website with \texttt{BEA\_API\_KEY})

The BEA provides detailed data on US-headquartered majority-owned foreign
affiliates (MOFAs), including employment, assets, sales, and value added by
country and industry.  These data complement the OECD CbCR statistics by
providing a longer time series (back to the 1980s) for the US parent population.

\textbf{Target variables (once downloaded):}
\begin{itemize}[noitemsep]
    \item Employment by country and industry (\texttt{bea\_mofa\_employment.xls})
    \item Total assets by country and industry (\texttt{bea\_mofa\_assets.xls})
    \item Value added by country and industry (\texttt{bea\_mofa\_value\_added.xls})
    \item Selected country-level aggregates (\texttt{bea\_mofa\_selected\_by\_country.xls})
\end{itemize}

\subsection{Eurostat Inward Foreign Affiliate Statistics (FATS)}
\label{subsec:eurostat}

\textbf{Canonical file:} \texttt{raw/eurostat\_fats/fats\_inward\_eu\_full.csv.gz}\\
\textbf{Supporting raw files:} \texttt{raw/eurostat\_fats/fats\_g1b\_08.csv.gz};
\texttt{raw/eurostat\_fats/fats\_activ.csv.gz}\\
\textbf{Eurostat datasets:} \texttt{fats\_g1b\_08} (2008--2020) and
\texttt{fats\_activ} (2021 onward)\\
\textbf{Combined rows:} 4{,}623{,}851 \quad \textbf{Combined size:} 21.44~MB\\
\textbf{Years:} 2008--2023 \quad \textbf{Status:} Downloaded and rebuilt from the
live Eurostat API on 10~March~2026

Eurostat inward FATS provide information on foreign-controlled enterprises
operating inside European reporting economies.  The data are useful for
descriptive validation, sectoral context, and external checks on the location of
real activity.  They are not a substitute for the OECD CbCR data, because they
do not contain tax liabilities or effective tax rates.

The Eurostat inward FATS series required a documented rebuild.  The repository
initially contained only \texttt{fats\_g1b\_08.csv.gz}, and inspection showed
that this file ended in 2020.  The live Eurostat metadata showed that Eurostat
had split the public inward FATS series into two official datasets: the legacy
series \texttt{fats\_g1b\_08} for 2008--2020 and the newer series
\texttt{fats\_activ} for 2021 onward.  The collection script
\texttt{data/collect\_all\_data.py} was therefore updated to download both
datasets through the Eurostat SDMX 2.1 API
(\url{https://ec.europa.eu/eurostat/api/dissemination/sdmx/2.1/data/}),
preserve both raw files, and then build one canonical project file.

\noindent\textbf{Collection and merge procedure:}
\begin{enumerate}[noitemsep]
    \item The existing file \texttt{fats\_g1b\_08.csv.gz} was inspected and
    confirmed to cover 2008--2020 only.
    \item The live Eurostat API metadata were checked to identify the current
    continuation series.  Eurostat labels \texttt{fats\_g1b\_08} as
    \emph{Foreign control of enterprises by controlling countries (2008--2020)}
    and \texttt{fats\_activ} as \emph{Foreign controlled EU enterprises by NACE
    Rev.2 activity and country of control (from 2021 onwards)}.
    \item The collection script was changed to download both source files from
    the API and to store them unchanged in \texttt{raw/eurostat\_fats/}.
    \item A canonical file, \texttt{fats\_inward\_eu\_full.csv.gz}, was built by
    vertically appending the two non-overlapping year ranges.
    \item Provenance columns were added to the canonical file:
    \texttt{source\_dataset}, \texttt{indicator\_scheme}, and
    \texttt{indicator\_code}.  This was done so that the series break remains
    visible in any later analysis.
    \item The rebuilt file was validated locally.  The final combined file
    covers 2008--2023 and fully covers the OECD CbCR window of 2016--2022.
\end{enumerate}

\noindent\textbf{Files created and verified:}
\begin{itemize}[noitemsep]
    \item \texttt{fats\_g1b\_08.csv.gz}: 1{,}421{,}600 rows; 6.22~MB; years 2008--2020
    \item \texttt{fats\_activ.csv.gz}: 3{,}202{,}251 rows; 15.91~MB; years 2021--2023
    \item \texttt{fats\_inward\_eu\_full.csv.gz}: 4{,}623{,}851 rows; 21.44~MB; years 2008--2023
\end{itemize}

The merge was intentionally conservative.  The two source files do not overlap
in years, so the merge is a stacked append and not a reconciliation of
conflicting observations.  That part is clean.  The difficult point lies in the
change in Eurostat's published structure from 2021 onward.

\noindent\textbf{Points that require attention in the merge:}
\begin{itemize}[noitemsep]
    \item The 2021 transition is a \textbf{series break}.  It should not be
    treated as a seamless continuation without explicit harmonisation.
    \item The activity dimension changes materially.  The old file has one broad
    aggregate, \texttt{B-N\_S95\_X\_K}, while the new file has 48 NACE-based
    activity categories plus a different top aggregate,
    \texttt{B-S\_X\_O\_S94}.
    \item The indicator system changes materially.  The old file uses 23 legacy
    \texttt{indic\_sb} codes such as \texttt{V12110} and \texttt{V12150}, while
    the new file uses 19 \texttt{indic\_sbs} codes such as
    \texttt{NETTUR\_MEUR} and \texttt{AV\_MEUR}.
    \item The country-of-control categories change materially.  The old file
    contains 255 control-country categories, while the new file contains 50.
    \item Several concepts are close but not obviously identical across the
    break.  For example, turnover in the old file becomes net turnover in the
    new file, personnel costs become employee benefits expense, and the
    investment concept is renamed and narrowed to tangible non-current assets.
    \item The canonical file therefore preserves the original scheme names and
    does not impose a forced one-to-one crosswalk at collection stage.
\end{itemize}

\noindent\textbf{Practical implications for the thesis methodology:}
\begin{itemize}[noitemsep]
    \item Eurostat FATS should be used as a supporting EU dataset.  Its strongest
    role is descriptive validation, sectoral context, and plausibility checks on
    the location of real activity.
    \item Eurostat FATS should not be treated as the core empirical base for the
    GMT estimates.  The OECD CbCR data remain the core source because they
    contain the profit and tax variables needed for the GMT simulations.
    \item Any figure or table that spans both sides of the 2021 break should
    either restrict attention to clearly harmonised concepts or mark the break
    explicitly.
    \item A pre/post design that uses the full 2008--2023 FATS window requires a
    separate harmonisation step.  That step should map only clearly comparable
    indicators, restrict to the closest common activity aggregate, and include
    an explicit post-2021 break indicator.
    \item For the current thesis design, the safest use is to draw on the
    post-2021 \texttt{fats\_activ} structure for recent EU sectoral evidence and
    to use the longer combined file mainly for traceability and broad descriptive
    background.
\end{itemize}


\section{Summary Table of All Datasets}
\label{sec:summary}

\begin{landscape}
\begin{longtable}{L{3.0cm} L{3.2cm} C{1.4cm} C{1.4cm} C{1.6cm} C{1.3cm} L{4.0cm}}
    \caption{Summary of all collected datasets}
    \label{tab:summary}\\
    \toprule
    \textbf{Dataset} & \textbf{File} & \textbf{Years} & \textbf{Countries} & \textbf{Rows} & \textbf{Size (MB)} & \textbf{Key variables} \\
    \midrule
    \endfirsthead
    \multicolumn{7}{c}{\textit{(continued from previous page)}} \\
    \toprule
    \textbf{Dataset} & \textbf{File} & \textbf{Years} & \textbf{Countries} & \textbf{Rows} & \textbf{Size} & \textbf{Key variables} \\
    \midrule
    \endhead
    \midrule
    \multicolumn{7}{r}{\textit{(continued on next page)}} \\
    \endfoot
    \bottomrule
    \endlastfoot
    %
    OECD CbCR Table II & \texttt{oecd\_cbcr/ cbcr\_table2\ldots} & 2016--2022 & 46 & 437{,}129 & 43.4 & PROFIT, TAX\_ACCRUED, TAX\_PAID, EMPLOYEES, TOT\_REV, ASSETS \\[4pt]
    OECD CbCR Table IV & \texttt{oecd\_cbcr/ cbcr\_table4\ldots} & 2016--2022 & 49 & 101{,}399 & 10.1 & Same + 13 activity types (ACT\_*) \\[4pt]
    OECD CbCR Table V & \texttt{oecd\_cbcr/ cbcr\_table5\ldots} & 2016--2022 & 40 & 4{,}195 & 0.4 & ASSETS, EMPLOYEES, UPR (P10/P90) \\[4pt]
    OECD CbCR Table VI & \texttt{oecd\_cbcr/ cbcr\_table6\ldots} & 2022 & 39 & 105{,}753 & 10.7 & Full measures + activity by MNE size \\[4pt]
    \midrule
    Statutory CIT rates & \texttt{oecd\_tax\_rates/ oecd\_statutory\ldots} & 2000--2025 & 144 & 13{,}273 & 1.2 & CIT (combined), CIT\_C (central), CIT\_SCG (sub-central) \\[4pt]
    Combined dividend rates & \texttt{oecd\_tax\_rates/ oecd\_combined\ldots} & 2000--2025 & 38 & 7{,}915 & 0.8 & CITPITCIT, FWHT, IR, IDTC \\[4pt]
    \midrule
    ETR baseline & \texttt{oecd\_etr/ oecd\_etr\_baseline} & 2017--2024 & 103 & 29{,}489 & 3.0 & EATR, EMTR, PBAR, CAPITAL\_ALLOWANCES \\[4pt]
    ETR income incentives & \texttt{oecd\_etr/ oecd\_etr\_income\ldots} & 2000--2024 & 55 & 5{,}784 & 0.7 & EATR\_ENH, PBAR\_ENH (IP~box) \\[4pt]
    ETR expenditure incentives & \texttt{oecd\_etr/ oecd\_etr\_expend\ldots} & 2019--2024 & 51 & 1{,}224 & 0.1 & EATR\_ENH, PBAR\_ENH (R\&D credits) \\[4pt]
    \midrule
    OECD avg.\ annual wages & \texttt{oecd\_wages/ oecd\_avg\ldots} & 1990--2024 & 41 & 3{,}924 & 0.3 & WG (USD PPP, national currency) \\[4pt]
    OECD hourly earnings & \texttt{oecd\_wages/ oecd\_hourly\ldots} & 1955--2026 & 41 & 41{,}421 & 3.4 & EAR (index, 2015=100) \\[4pt]
    ILO hourly earnings & \texttt{ilo\_wages/ ilo\_hourly\ldots} & 1990--2025 & 105 & 227{,}317 & 50.4 & EAR\_EHRA\_NB (LCU, PPP) \\[4pt]
    \midrule
    WB GDP per capita & \texttt{world\_bank/ wb\_gdp\ldots} & 2015--2024 & 217 & $\sim$2{,}170 & $<$0.1 & gdp\_pc\_usd \\[4pt]
    WB population & \texttt{world\_bank/ wb\_population} & 2015--2024 & 217 & $\sim$2{,}170 & $<$0.1 & population \\[4pt]
    WB income groups & \texttt{world\_bank/ wb\_income\ldots} & Cross-section & 299 & 299 & $<$0.1 & income\_level, income\_id, region \\[4pt]
    \midrule
    CbCR requirements & \texttt{pillar2\_beps/ oecd\_cbcr\_req\ldots} & Cross-section & 119 & 493 & $<$0.1 & CBCR\_LAW, THRESHOLD, DEADLINE \\[4pt]
    CFC rules & \texttt{pillar2\_beps/ oecd\_cfc\ldots} & 2019,24,25 & 146 & 1{,}406 & 0.4 & CFC, TRIGGER\_RATE, YEAR \\[4pt]
    Hybrid mismatch & \texttt{pillar2\_beps/ oecd\_hybrid\ldots} & 2025 & 145 & 985 & 0.1 & ACT2, REC1--REC15, YEAR \\[4pt]
    Interest limitation & \texttt{pillar2\_beps/ oecd\_interest\ldots} & 2019,24,25 & 146 & 3{,}633 & 0.3 & ILR, GR, TR, FAM, CF \\[4pt]
    IP regimes & \texttt{pillar2\_beps/ oecd\_ip\ldots} & 2018--2025 & 64 & 2{,}282 & 0.3 & NAME, RATE, NEXUS, YEAR\_START \\[4pt]
    Mandatory disclosure & \texttt{pillar2\_beps/ oecd\_mandatory\ldots} & 2025 & 145 & 985 & 0.1 & ACT12, YEAR \\[4pt]
    Pillar Two tracker & \texttt{pillar2\_qdmtt/ pwc\_pillar2\ldots} & multi-date & multiple & --- & $<$0.1 & IIR, UTPR, QDMTT status and dates \\[4pt]
    Eurostat inward FATS (combined) & \texttt{eurostat\_fats/ fats\_inward\_eu\_full.csv.gz} & 2008--2023 & 32--33 & 4{,}623{,}851 & 21.4 & \texttt{indicator\_code}, \texttt{c\_ctrl}, \texttt{geo}, \texttt{nace\_r2}, \texttt{obs\_value} \\[4pt]
\end{longtable}
\end{landscape}


\section{Data Limitations and Methodological Notes}
\label{sec:limitations}

\begin{enumerate}

    \item \textbf{CbCR aggregation bias.}  CbCR data are aggregated by the OECD
    to prevent identification of individual companies.  Interquartile means
    (Table~II) reduce but do not eliminate outlier influence; extreme values in
    Table~IV suggest residual sensitivity.  Researchers should apply winsorising
    or log transformations when using level variables.

    \item \textbf{CbCR reporting gaps.}  Not all MNE home countries have enacted
    CbCR legislation.  Key gaps include the United States (which filed with the
    OECD only from 2020), Russia (no OECD membership), and most developing
    countries.  The bilateral coverage matrix (REF\_AREA $\times$
    COUNTERPART\_AREA) is therefore unbalanced.

    \item \textbf{ETR model assumptions.}  The OECD ETR calculations assume a
    representative investment (equal weights on machinery, buildings, and
    intangibles) and a fixed nominal pre-tax rate of return.  They do not capture
    firm-specific or industry-specific differences.  The choice of
    \texttt{ETR\_TAX\_TYPE} (consolidated vs.\ fixed-rate) materially affects
    results and should be justified.

    \item \textbf{Statutory vs.\ effective rates.}  Statutory CIT rates
    (Section~\ref{sec:cit}) and ETRs (Section~\ref{sec:etr}) measure different
    concepts.  The statutory rate is the relevant policy instrument; the EATR
    captures actual investment incentives including depreciation and incentives.
    Both may diverge from the \emph{backward-looking} ETRs implicit in the CbCR
    (TAX\_ACCRUED / PROFIT).

    \item \textbf{Wage data comparability.}  OECD wages (USD PPP, 2024 base)
    are directly comparable across countries.  ILO wages are in local currencies
    and require PPP conversion for cross-country analysis.  The ILO
    \texttt{CUR\_TYPE\_PPP} observations provide a consistent alternative.

    \item \textbf{BEPS policy indicators.}  The cross-sectional BEPS datasets
    (CbCR requirements, hybrid mismatch, mandatory disclosure) have no
    time dimension; they reflect the current status only.  CFC rules and
    interest limitation data are available for 2019, 2024, and 2025, providing
    a limited before/after comparison.  For richer difference-in-differences
    designs, the specific year of enactment (\texttt{YEAR} variable) should be
    used to construct country $\times$ time indicators.

    \item \textbf{Pillar Two tracker currency.}  The PwC tracker is maintained on
    a commercial basis and may contain inaccuracies or lag behind the most
    recent legislative changes.  Cross-referencing against the OECD Central
    Record of Qualified Domestic Minimum Taxes is advisable for publications.

    \item \textbf{BEA availability.}  The BEA files still require a free
    \texttt{BEA\_API\_KEY} and were not downloaded in the present iteration.

    \item \textbf{Eurostat FATS series break.}  Eurostat inward FATS are now
    available in the repository, but the public series changes in 2021.  The
    pre-2021 file (\texttt{fats\_g1b\_08}) and the post-2020 file
    (\texttt{fats\_activ}) differ in activity classification, indicator scheme,
    and country-of-control structure.  The canonical file therefore preserves
    provenance and should not be interpreted as a fully harmonised panel.

    \item \textbf{Eurostat FATS use in the methodology.}  For this thesis,
    Eurostat FATS are best treated as a supporting EU dataset.  They can
    strengthen descriptive claims and sectoral validation, but any pre/post or
    long-run econometric use must first justify a transparent harmonisation
    strategy around the 2021 break.

\end{enumerate}

\end{document}
